% problem_id: S046-thm:orth-fact
% source_id: S046 (HoTT Book)
% source_file: hlevels.tex
% statement_locator: hlevels.tex lines 1511--1524
% proof_locator: hlevels.tex lines 1526--1573
% env: thm
% pairing: tight (verified: the next environment after the statement is its proof)
% status: complete (statement + author proof verbatim + reason + locators)
%
\begin{thm}\label{thm:orth-fact}
For each $f:A\to B$, the space $\fact_n(f)$ defined by
\begin{equation*}
\sm{X:\type}{g:A\to X}{h:X\to B} (h\circ g\htpy f)\times\mathsf{conn}_n(g)\times\mathsf{trunc}_n(h)
\end{equation*}
is contractible.
Its center of contraction is the element
\begin{equation*}
\pairr{\im_n(f),\tilde{f},\proj1,\theta,\varphi,\psi}:\fact_n(f)
\end{equation*}
arising from \cref{prop:to_image_is_connected},
where $\theta:\proj1\circ\tilde{f}\htpy f$ is the canonical homotopy, where $\varphi$ is the proof of
\cref{prop:to_image_is_connected}, and where $\psi$ is the obvious proof that $\proj1:\im_n(f)\to B$ has $n$-truncated fibers.
\end{thm}

\begin{proof}
By \cref{prop:to_image_is_connected} we know that there is an element of $\fact_n(f)$, hence it is enough to
show that $\fact_n(f)$ is a mere proposition. Suppose we have two $n$-factorizations
\begin{equation*}
\pairr{X_1,g_1,h_1,H_1,\varphi_1,\psi_1}\qquad\text{and}\qquad\pairr{X_2,g_2,h_2,H_2,\varphi_2,\psi_2}
\end{equation*}
of $f$. Then we have the pointwise-concatenated homotopy
\[ H\defeq (\lam{a} H_1(a) \ct H_2^{-1}(a)) \,:\, (h_1\circ g_1\htpy h_2\circ g_2).\]
By univalence and the characterization of paths and transport in $\Sigma$-types, function types, and path types, it suffices to show that
\begin{enumerate}
\item there is an equivalence $e:X_1\eqvsym X_2$,
\item there is a homotopy $\zeta:e\circ g_1\htpy g_2$,
% \note{Is it easy enough to see that these elements are the various transports?}
\item there is a homotopy $\eta:h_2\circ e\htpy h_1$,
\item for any $a:A$ we have $\opp{\apfunc{h_2}(\zeta(a))} \ct \eta(g_1(a)) \ct H_1(a) = H_2(a)$.
\end{enumerate}
%where $\underline{e}$ is the function underlying the equivalence.
We prove these four assertions in that order.
\begin{enumerate}
\item By \cref{prop:factor_equiv_fiber}, we have a fiberwise equivalence
% \note{It could be a nice exercise for the book to show
% that if $f_1:A_1\to B$ and $f_2:A_2\to B$ have equivalent fibers, then $A_1\eqvsym A_2$}.
\begin{equation*}
E(H) : \prd{b:B} \eqv{\hfib{h_1}b}{\hfib{h_2}b}.
\end{equation*}
This induces an equivalence of total spaces, i.e.\ we have
\begin{equation*}
\eqvspaced{\Parens{\sm{b:B} \hfib{h_1}b}}{\Parens{\sm{b:B}\hfib{h_2}b}}.
\end{equation*}
Of course, we also have the equivalences $X_1\eqvsym\sm{b:B}\hfib{h_1}b$ and $X_2\eqvsym\sm{b:B}
\hfib{h_2}b$ from \cref{thm:total-space-of-the-fibers}.
This gives us our equivalence $e:X_1\eqvsym X_2$; the reader may verify that the underlying function of $e$ is given by
\begin{equation*}
e(x) \jdeq \proj1(E(H,h_1(x))(x,\refl{h_1(x)})).
\end{equation*}
\item By \cref{prop:factor_equiv_fiber}, we may choose
  $\zeta(a) \defeq \apfunc{\proj1}(\overline E(H,a)) : e(g_1(a)) = g_2(a)$.
  \label{item:orth-fact-2}
\item For every $x:X_1$, we have
\begin{equation*}
\proj2(E(H,h_1(x))({x,\refl{h_1(x)}})) :h_2(e(x))= h_1(x),
\end{equation*}
giving us a homotopy $\eta:h_2\circ e\htpy h_1$.
\item By the characterization of paths in fibers (\cref{lem:hfib}), the path $\overline E(H,a)$ from \cref{prop:factor_equiv_fiber} gives us
  $\eta(g_1(a)) = \apfunc{h_2}(\zeta(a)) \ct \opp{H(a)}$.
  The desired equality follows by substituting the definition of $H$ and rearranging paths.\qedhere
\end{enumerate}
\end{proof}
